\section{Edit zoo}
The builder implements six families. Random-direction ablation projects selected weight matrices away from sampled directions. Projected abliteration uses refusal directions computed as differences of mean activations on harmful and harmless prompts. Its Heretic-parity classic kernel applies $W'=W-w_\ell v(v^\top W)$ with no row normalization. The layer weight $w_\ell$ interpolates linearly from a peak to a floor within a finite radius and is zero outside it. Directions may be selected or interpolated across residual layers. This is a specific projected variant, not Heretic's default row normalization.

Heretic ARA generates an adapter using Optuna trials and refusal and KL scores; a configured KL constraint restricts eligible trials. Heretic describes ARA as its default modifier and uses Optuna TPE optimization \citep{heretic2026}. LoRA supervised fine-tuning trains adapters on GSM8K or Alpaca-style data with sampled training parameters. Task-vector entries scale an eligible parent delta relative to the pinned base. Chains store cumulative snapshots relative to the base. Random-direction and abliteration steps modify the accumulated state, whereas adapter steps are optimized against the base and added to that state. Consecutive snapshots must be differenced to obtain step deltas; a snapshot must not be applied again to its parent.

An entry ID hashes the family, parameters, seed, pinned base revision and parent. A separate reuse key records further provenance, including kernel, dtype, refusal fingerprint and chain prefix; mismatches prevent silent reuse. Eligibility distinguishes production, smoke and legacy unclassified entries. Smoke status propagates through descendants, and task-vector parents must have compatible base and provenance and normally be production entries.

Table~\ref{tab:zoo-counts} counts manifest rows in the local build at the time of writing, excluding manifest headers. Chain rows count cumulative steps rather than independent trajectories. The manifest is a build artefact, not evidence that every row has a completed evaluation.
\begin{table}[t]
\centering
\begin{tabular}{lrrr}\toprule
Family & Production & Smoke & Unclassified \\\midrule
Random direction & 150 & 0 & 3 \\
Abliteration & 200 & 0 & 5 \\
Heretic ARA & 0 & 0 & 1 \\
LoRA SFT & 120 & 0 & 1 \\
Task vector & 80 & 0 & 0 \\
Chain steps & 523 & 0 & 20 \\\bottomrule
\end{tabular}
\caption{Local \texttt{data/zoo/qwen3.5-0.8b/manifest.jsonl}, read 4 November 2026. Counts are rows by recorded eligibility; unclassified denotes legacy rows without that field.}
\label{tab:zoo-counts}
\end{table}
